\chapter{Matrices y grafos}\label{ch:g12:matrix}

Una \omterm{def:g12:matrix:matrix}{matriz} es una tabla rectangular de números que se suman y se
multiplican según reglas diseñadas para que el álgebra de matrices
represente la composición de transformaciones lineales. Las matrices
resuelven \omterm{def:g10:lines:system}{sistemas lineales}, gobiernan sucesiones recurrentes acopladas y
cuentan caminos en redes: son las matemáticas que hay detrás de los
buscadores de internet y de los algoritmos de camino más corto.

\section{Álgebra de matrices}

\begin{definition}[Matriz]\label{def:g12:matrix:matrix}
Una \emph{matriz $m \times n$}\index{matriz} es una tabla de \omterm{def:g10:numbers:sets}{números reales}
con $m$ filas y $n$ columnas: $A = (a_{ij})$, donde $a_{ij}$ es el
coeficiente de la fila $i$ y la columna $j$. Dos matrices del mismo tamaño
se suman coeficiente a coeficiente, y $\lambda A = (\lambda a_{ij})$.
\end{definition}

\begin{definition}[Producto de matrices]\label{def:g12:matrix:product}
Sea $A$ de tamaño $m \times n$ y sea $B$ de tamaño $n \times p$. El
producto $AB$ es la \omterm{def:g12:matrix:matrix}{matriz} $m \times p$ cuyo coeficiente $(i,j)$ es
\[
(AB)_{ij} = \sum_{k=1}^{n} a_{ik} b_{kj}
\]
(la regla ``fila $i$ de $A$ por columna $j$ de $B$'').
\end{definition}

\begin{example}
$\begin{pmatrix} 1 & 2\\ 3 & 4\end{pmatrix}
\begin{pmatrix} 0 & 1\\ 1 & 1\end{pmatrix}
= \begin{pmatrix} 2 & 3\\ 4 & 7\end{pmatrix}$,
\quad mientras que \quad
$\begin{pmatrix} 0 & 1\\ 1 & 1\end{pmatrix}
\begin{pmatrix} 1 & 2\\ 3 & 4\end{pmatrix}
= \begin{pmatrix} 3 & 4\\ 4 & 6\end{pmatrix}$:
\emph{el producto de matrices no es conmutativo}.
\end{example}

\begin{proposition}[Reglas del álgebra de matrices]\label{prop:g12:matrix:rules}
Siempre que los tamaños den sentido a los productos:
\[
(AB)C = A(BC), \qquad A(B + C) = AB + AC, \qquad (A+B)C = AC + BC,
\]
y la \emph{\omterm{def:g12:matrix:matrix}{matriz} identidad} $I_n$ (unos en la diagonal y ceros fuera de
ella) cumple $I_m A = A I_n = A$ para $A$ de tamaño $m \times n$.
\end{proposition}

\begin{proof}
Todas son comprobaciones coeficiente a coeficiente a partir de la
\cref{def:g12:matrix:product}; la asociatividad, la única no trivial, se
reduce a intercambiar dos sumas finitas:
$\bigl((AB)C\bigr)_{ij} = \sum_l \left(\sum_k a_{ik}b_{kl}\right) c_{lj}
= \sum_k a_{ik} \left(\sum_l b_{kl} c_{lj}\right)
= \bigl(A(BC)\bigr)_{ij}$.
\end{proof}

\begin{definition}[Inversa]\label{def:g12:matrix:inverse}
Una \omterm{def:g12:matrix:matrix}{matriz} cuadrada $A$ de tamaño $n$ es
\emph{invertible}\index{matriz!inversa} si existe una \omterm{def:g12:matrix:matrix}{matriz} $B$ con
$AB = BA = I_n$; esa $B$ es entonces única y se escribe $A^{-1}$.
\end{definition}

\begin{proposition}[Inversa de una matriz $2\times2$]\label{prop:g12:matrix:inverse2x2}
Sea $A = \begin{pmatrix} a & b\\ c & d\end{pmatrix}$ y sea
$\det A = ad - bc$ (el \emph{determinante}\index{determinante}). Entonces
$A$ es \omterm{def:g12:matrix:inverse}{invertible} si y solo si $\det A \neq 0$, y en ese caso
\[
A^{-1} = \frac{1}{ad - bc}\begin{pmatrix} d & -b\\ -c & a\end{pmatrix}.
\]
\end{proposition}

\begin{proof}
Un cálculo da
$A \begin{pmatrix} d & -b\\ -c & a\end{pmatrix}
= \begin{pmatrix} d & -b\\ -c & a\end{pmatrix} A = (ad - bc) I_2$; si
$ad - bc \neq 0$, dividimos. Recíprocamente, si $ad - bc = 0$, las columnas
de $A$ son proporcionales, y también lo son las de $AB$ sea cual sea $B$;
pero las columnas de $I_2$ no son proporcionales, así que ninguna $B$ puede
cumplir $AB = I_2$.
\end{proof}

\begin{method}[Sistemas lineales]\label{met:g12:matrix:systems}
El sistema
$\begin{cases} ax + by = e\\ cx + dy = f \end{cases}$
es la \omterm{def:g10:algebra:equation}{ecuación} matricial $AX = Y$ con
$X = \begin{pmatrix} x \\ y\end{pmatrix}$ e
$Y = \begin{pmatrix} e \\ f\end{pmatrix}$. Si $\det A \neq 0$, su única
solución es $X = A^{-1}Y$. El mismo formalismo sirve para $n$ ecuaciones
con $n$ incógnitas.
\end{method}

\section{Potencias de matrices y sucesiones recurrentes}

\begin{definition}\label{def:g12:matrix:power}
Para una \omterm{def:g12:matrix:matrix}{matriz} cuadrada $A$ y $k \in \N$, $A^k = A \times \dots \times A$
($k$ factores), con $A^0 = I$.
\end{definition}

\begin{method}[Casos diagonal más nilpotente y diagonalizable]\label{met:g12:matrix:powers}
Dos maneras estándar de calcular $A^k$:
\begin{itemize}
  \item Si $A = \lambda I + N$ con $N^2 = 0$, el teorema del binomio
        (válido aquí porque $I$ y $N$ conmutan) se reduce a dos términos:
        $A^k = \lambda^k I + k \lambda^{k-1} N$.
  \item Si se encuentra una \omterm{def:g12:matrix:matrix}{matriz} \omterm{def:g12:matrix:inverse}{invertible} $P$ con $A = PDP^{-1}$ y $D$
        diagonal, entonces $A^k = P D^k P^{-1}$, y $D^k$ se calcula
        coeficiente a coeficiente. (Hallar tal $P$ de manera sistemática es
        la teoría de la \emph{diagonalización}, que se desarrolla en la
        universidad; a este nivel, $P$ viene dada.)
\end{itemize}
\end{method}

\begin{example}[Sucesiones acopladas]\label{ex:g12:matrix:coupled}
Sean $u_{n+1} = 3u_n + v_n$ y $v_{n+1} = u_n + 3v_n$. Tomando
$X_n = \begin{pmatrix} u_n\\ v_n \end{pmatrix}$ y
$A = \begin{pmatrix} 3 & 1\\ 1 & 3\end{pmatrix}$, obtenemos
$X_{n+1} = AX_n$, luego $X_n = A^n X_0$. Las sucesiones auxiliares
$s_n = u_n + v_n$ y $d_n = u_n - v_n$ cumplen $s_{n+1} = 4s_n$ y
$d_{n+1} = 2d_n$, así que $s_n = 4^n s_0$, $d_n = 2^n d_0$ y
\[
u_n = \frac{4^n(u_0+v_0) + 2^n(u_0-v_0)}{2}, \qquad
v_n = \frac{4^n(u_0+v_0) - 2^n(u_0-v_0)}{2}.
\]
(Entre bastidores: $(1,1)$ y $(1,-1)$ son direcciones propias de $A$.)
\end{example}

\section{Grafos y caminos}

\begin{definition}[Grafo, matriz de adyacencia]\label{def:g12:matrix:graph}
Un \emph{grafo}\index{grafo} consta de vértices $1, 2, \dots, n$ y de
aristas que unen ciertos pares de vértices (pares ordenados si el grafo es
\emph{dirigido}). Su \emph{matriz de adyacencia}\index{matriz de
adyacencia} es la \omterm{def:g12:matrix:matrix}{matriz} $n \times n$ $M$ con $m_{ij} = 1$ si hay una
arista de $i$ a $j$, y $0$ en caso contrario. Un \emph{camino} de longitud
$k$ de $i$ a $j$ es una \omterm{def:g12:seq:sequence}{sucesión} de $k$ aristas consecutivas que lleva de
$i$ a $j$.
\end{definition}

\begin{omfigure}
\begin{tikzpicture}[baseline=(current bounding box.center)]
  \node[circle, draw, inner sep=2.5pt] (v1) at (0,0) {$1$};
  \node[circle, draw, inner sep=2.5pt] (v2) at (2.6,1.2) {$2$};
  \node[circle, draw, inner sep=2.5pt] (v3) at (2.6,-1.2) {$3$};
  \draw[-Stealth, omDef, thick] (v1) -- (v2);
  \draw[-Stealth, omDef, thick] (v2) -- (v3);
  \draw[-Stealth, omDef, thick] (v1) to[bend right=18] (v3);
  \draw[-Stealth, omDef, thick] (v3) to[bend right=18] (v1);
\end{tikzpicture}
\qquad\qquad
$M = \begin{pmatrix} 0 & 1 & 1\\ 0 & 0 & 1\\ 1 & 0 & 0 \end{pmatrix}$

{\small Un \omterm{def:g12:matrix:graph}{grafo} dirigido y su \omterm{def:g12:matrix:graph}{matriz de adyacencia}
(\cref{exo:g12:matrix:6}): $m_{ij} = 1$ exactamente cuando hay una arista
de $i$ a $j$.}
\end{omfigure}

\begin{theorem}[Recuento de caminos]\label{thm:g12:matrix:walks}
El número de caminos de longitud $k$ del vértice $i$ al vértice $j$ es el
coeficiente $(i,j)$ de $M^k$.
\end{theorem}

\begin{proof}
Inducción sobre $k$. Para $k = 1$ es la definición de $M$. Supongamos el
resultado para $k$. Un camino de longitud $k+1$ de $i$ a $j$ es un camino
de longitud $k$ de $i$ a cierto vértice $l$, seguido de una arista de $l$ a
$j$; por los principios de la suma y del producto, su número es
\[
\sum_{l=1}^{n} \bigl(M^k\bigr)_{il}\, m_{lj} = \bigl(M^{k+1}\bigr)_{ij}.
\qedhere
\]
\end{proof}

\begin{example}
Para el \omterm{def:g12:matrix:graph}{grafo} triángulo ($3$ vértices, todos los pares unidos),
$M = \begin{pmatrix} 0&1&1\\ 1&0&1\\ 1&1&0\end{pmatrix}$ y
$M^2 = \begin{pmatrix} 2&1&1\\ 1&2&1\\ 1&1&2\end{pmatrix}$: desde cada
vértice hay $2$ caminos de longitud $2$ que vuelven a él (por uno u otro
vecino) y $1$ hacia cada uno de los otros vértices.
\end{example}

\section{Ejercicios}

\begin{exercise}[$\star$]\label{exo:g12:matrix:1}
Sean $A = \begin{pmatrix} 1 & 2\\ 0 & 1 \end{pmatrix}$ y
$B = \begin{pmatrix} 2 & 0\\ 1 & 1 \end{pmatrix}$. Calcula $A + B$, $AB$,
$BA$ y $A^2$.
\end{exercise}

\begin{exercise}[$\star$]\label{exo:g12:matrix:2}
Determina si las matrices siguientes son \omterm{def:g12:matrix:inverse}{invertibles} y calcula las inversas
cuando existan:
\[
A = \begin{pmatrix} 2 & 5\\ 1 & 3\end{pmatrix}, \qquad
B = \begin{pmatrix} 3 & 6\\ 2 & 4\end{pmatrix}.
\]
\end{exercise}

\begin{exercise}[$\star$]\label{exo:g12:matrix:3}
Resuelve por inversión matricial el sistema
$\begin{cases} 2x + 5y = 1\\ x + 3y = 2 . \end{cases}$
\end{exercise}

\begin{exercise}[$\star\star$]\label{exo:g12:matrix:4}
Sea $A = \begin{pmatrix} 2 & 1\\ 0 & 2\end{pmatrix} = 2I + N$ con
$N = \begin{pmatrix} 0 & 1\\ 0 & 0 \end{pmatrix}$.
\begin{enumerate}
  \item Comprueba que $N^2 = 0$ y que $I$ y $N$ conmutan.
  \item Deduce $A^k$ para todo $k \in \N$ y verifica la fórmula para $k=2$
        mediante un cálculo directo.
\end{enumerate}
\end{exercise}

\begin{exercise}[$\star\star$]\label{exo:g12:matrix:5}
Sea $A = \begin{pmatrix} 0 & 1\\ 1 & 1\end{pmatrix}$ y sea $F_n$ la
\omterm{def:g12:seq:sequence}{sucesión} de Fibonacci ($F_0 = 0$, $F_1 = 1$, $F_{n+2} = F_{n+1} + F_n$).
Demuestra por inducción que, para $n \geq 1$,
\[
A^n = \begin{pmatrix} F_{n-1} & F_n\\ F_n & F_{n+1}\end{pmatrix},
\]
y deduce la identidad $F_{n+1}F_{n-1} - F_n^2 = (-1)^n$. (Indicación: los
\omterm{def:g11:vect:det}{determinantes} se multiplican: $\det(MN) = \det M \det N$, cosa que puedes
comprobar para matrices $2\times2$.)
\end{exercise}

\begin{exercise}[$\star\star$]\label{exo:g12:matrix:6}
Un \omterm{def:g12:matrix:graph}{grafo} dirigido sobre los vértices $\{1, 2, 3\}$ tiene las aristas
$1\to2$, $2\to3$, $3\to1$ y $1\to3$.
\begin{enumerate}
  \item Escribe la \omterm{def:g12:matrix:graph}{matriz de adyacencia} $M$ y calcula $M^2$ y $M^3$.
  \item ¿Cuántos caminos de longitud $3$ van de $1$ a $1$? Enuméralos.
\end{enumerate}
\end{exercise}

\begin{exercise}[$\star\star$]\label{exo:g12:matrix:7}
Una empresa de coches compartidos mueve vehículos entre dos ciudades $A$ y
$B$. Cada semana, el $80\,\%$ de los coches que están en $A$ se quedan en
$A$ y el $20\,\%$ pasa a $B$; el $30\,\%$ de los coches que están en $B$
pasa a $A$ y el $70\,\%$ se queda. Sean $a_n$ y $b_n$ las proporciones de
la flota en cada ciudad.
\begin{omfigure}
\begin{tikzpicture}
  \node[circle, draw, inner sep=3pt] (A) at (0,0) {$A$};
  \node[circle, draw, inner sep=3pt] (B) at (3.2,0) {$B$};
  \draw[-Stealth, omDef, thick] (A) to[bend left=25]
    node[above, font=\small] {$0.2$} (B);
  \draw[-Stealth, omDef, thick] (B) to[bend left=25]
    node[below, font=\small] {$0.3$} (A);
  \draw[-Stealth, omDef, thick] (A) to[loop left]
    node[left, font=\small] {$0.8$} (A);
  \draw[-Stealth, omDef, thick] (B) to[loop right]
    node[right, font=\small] {$0.7$} (B);
\end{tikzpicture}
\end{omfigure}
\begin{enumerate}
  \item Escribe $X_{n+1} = MX_n$ con
        $X_n = \begin{pmatrix} a_n\\ b_n\end{pmatrix}$ e identifica $M$.
  \item Halla las proporciones de equilibrio (resuelve $MX = X$ con
        $a + b = 1$).
  \item Demuestra que $c_n = a_n - 0.6$ cumple $c_{n+1} = 0.5\,c_n$ y
        concluye que el reparto de la flota \omterm{def:g12:seq:limit}{converge} hacia el equilibrio.
\end{enumerate}
\end{exercise}

\begin{exercise}[$\star\star\star$]\label{exo:g12:matrix:8}
Sean $A = \begin{pmatrix} 3 & 1\\ 1 & 3 \end{pmatrix}$ y
$P = \begin{pmatrix} 1 & 1\\ 1 & -1 \end{pmatrix}$.
\begin{enumerate}
  \item Calcula $P^{-1}$ y después $D = P^{-1}AP$, y comprueba que $D$ es
        diagonal.
  \item Deduce una fórmula cerrada para $A^n$ y compárala con el
        \cref{ex:g12:matrix:coupled}.
\end{enumerate}
\end{exercise}

\section{Problema: La matriz que se sabe Fibonacci (y el tiempo que va a
hacer)}

\begin{problem}[{Problema de fin de semana --- una sola matriz $2 \times 2$
lleva dentro todo Fibonacci, una matriz de Markov predice el tiempo a largo
plazo y un vector propio vale mil millones de dólares}]
\label{pb:g12:matrix:1}
Una \omterm{def:g12:matrix:matrix}{matriz} es una máquina que se come un estado y devuelve el siguiente; y
sus \emph{potencias} guardan, por tanto, futuros enteros. Este problema
abre con la asombrosa \omterm{def:g12:matrix:matrix}{matriz} cuyas potencias enumeran los números de
Fibonacci (y demuestran sus identidades a línea por identidad), lleva
después el tiempo atmosférico, visto como cadena de Markov, hasta su estado
estacionario, y cierra con el \omterm{def:g11:vect:vector}{vector} propio sobre el que se construyó un
buscador (\cref{thm:g12:matrix:walks}, \cref{met:g12:matrix:powers}).

\medskip
\noindent\textbf{Parte I --- Soltura.}
\begin{enumerate}
  \item Con $A = \begin{pmatrix} 1 & 2\\ 3 & 4\end{pmatrix}$ y
        $B = \begin{pmatrix} 0 & 1\\ 1 & 0\end{pmatrix}$: calcula $AB$ y
        $BA$. ¿Veredicto sobre la conmutatividad?
  \item Invierte $\begin{pmatrix} 2 & 1\\ 5 & 3\end{pmatrix}$
        (\cref{prop:g12:matrix:inverse2x2}) y usa la inversa para resolver
        $2x + y = 4$, $5x + 3y = 7$.
  \item Sea $N = \begin{pmatrix} 0 & 1\\ 0 & 0\end{pmatrix}$: calcula
        $N^2$ y deduce que $(I + N)^n = I + nN$ para todo $n$.
  \item El \omterm{def:g12:matrix:graph}{grafo} triángulo (tres vértices, todos los pares unidos):
        escribe su \omterm{def:g12:matrix:graph}{matriz de adyacencia} $A$, calcula $A^3$ e interpreta los
        coeficientes diagonales (\cref{thm:g12:matrix:walks}).
  \item Para $D = \begin{pmatrix} 2 & 0\\ 0 & \frac12 \end{pmatrix}$: da
        $D^n$ y su comportamiento cuando $n \to \infty$.
\end{enumerate}

\medskip
\noindent\textbf{Parte II --- La \omterm{def:g12:matrix:matrix}{matriz} de Fibonacci.}
Sea $F = \begin{pmatrix} 1 & 1\\ 1 & 0\end{pmatrix}$ y sean
$F_1 = F_2 = 1, F_3 = 2, \dots$ los números de Fibonacci del
\cref{pb:g11:seq:1}.
\begin{enumerate}[resume]
  \item Calcula $F^2$, $F^3$ y $F^4$, y conjetura la forma general de $F^n$
        en términos de números de Fibonacci.
  \item Demuestra por inducción la conjetura
        $F^n = \begin{pmatrix} F_{n+1} & F_n\\ F_n & F_{n-1}
        \end{pmatrix}$.
  \item Toma \omterm{def:g11:vect:det}{determinantes} en los dos miembros (el \omterm{def:g11:vect:det}{determinante} de un
        producto es el producto de los \omterm{def:g11:vect:det}{determinantes}; compruébalo con
        matrices $2 \times 2$ si nunca lo has visto): deduce la
        \emph{identidad de Cassini}
        $F_{n+1}F_{n-1} - F_n^2 = (-1)^n$, el motor del cuadrado que se
        desvanece, demostrada en una línea.
  \item A partir de $F^{m+n} = F^m F^n$, lee los coeficientes superiores
        derechos y deduce la \emph{fórmula de adición}
        \[
        F_{m+n} = F_{m+1} F_n + F_m F_{n-1} .
        \]
        Compruébala para $m = n = 3$.
  \item Deduce de la fórmula de adición (por inducción sobre $k$) que $F_n$
        \omterm{def:g12:arith:divides}{divide} a $F_{kn}$, y verifícalo en $F_3 \mid F_6$ y
        $F_3 \mid F_9$.
  \item Para calcular $F_{100}$ no hace falta multiplicar $100$ matrices:
        basta elevar al cuadrado repetidamente ($F^2, F^4, F^8, \dots$) y
        combinar. ¿Cuántos productos de matrices bastan y de qué antiguo
        truco de multiplicación del volumen anterior se trata, ascendido a
        las matrices?
\end{enumerate}

\medskip
\noindent\textbf{Parte III --- La máquina del tiempo atmosférico.}
En cierta ciudad, tras un día soleado el siguiente es soleado con
\omterm{def:g10:proba:distribution}{probabilidad} $0.8$; y tras un día de lluvia, es soleado con \omterm{def:g10:proba:distribution}{probabilidad}
$0.4$. Codificamos la \omterm{def:g11:prob:rv}{distribución} del día como una columna
$\binom{p_{\text{sol}}}{p_{\text{lluvia}}}$ y la evolución mediante
\[
M = \begin{pmatrix} 0.8 & 0.4\\ 0.2 & 0.6 \end{pmatrix}.
\]
\begin{enumerate}[resume]
  \item Comprueba que cada columna de $M$ suma $1$ y explica por qué toda
        máquina del tiempo tiene que cumplir esa propiedad.
  \item Hoy hace sol. Calcula la previsión para mañana y para pasado
        mañana.
  \item Halla el \emph{estado estacionario}: la \omterm{def:g11:prob:rv}{distribución} $v$ con
        $Mv = v$ (y con coeficientes que sumen $1$). ¿Qué proporción de
        días es soleada a largo plazo?
  \item Parte de un día de lluvia, $\binom01$, y aplica $M$ cuatro veces,
        siguiendo en cada paso la distancia al estado estacionario. ¿Por
        qué factor se encoge la \omterm{def:g11:seq:arithmetic}{diferencia} en cada paso y de qué tipo de
        convergencia se trata?
  \item PageRank en miniatura: tres páginas, con enlaces $A \to B$,
        $A \to C$, $B \to C$ y $C \to A$. Un navegante aleatorio sigue un
        enlace saliente elegido al azar de manera uniforme. Escribe la
        \omterm{def:g12:matrix:matrix}{matriz} de transición, halla el estado estacionario y ordena las
        páginas.
  \item Interpreta la clasificación: ¿por qué puntúa $C$ tan alto como $A$
        pese a recibir enlaces de menos páginas? ¿Qué mide realmente el
        estado estacionario? (El PageRank de verdad añade un factor de
        amortiguación para los callejones sin salida y los saltos; la idea
        del \omterm{def:g11:vect:vector}{vector} propio es exactamente esta.)
\end{enumerate}

\medskip
\noindent\textbf{Parte IV --- Los dividendos de la diagonal.}
\begin{enumerate}[resume]
  \item Dos cantidades acopladas obedecen a $u_{n+1} = 3u_n + v_n$ y
        $v_{n+1} = u_n + 3v_n$, es decir, a la \omterm{def:g12:matrix:matrix}{matriz} $A$ del
        \cref{exo:g12:matrix:8}. Usando la diagonalización de ese ejercicio
        ($D = \operatorname{diag}(4, 2)$), da la fórmula cerrada de $u_n$
        cuando $u_0 = 1$ y $v_0 = 0$, y contrástala con el cálculo directo
        para $n = 1, 2, 3$.
  \item En una o dos frases: ¿qué le \emph{hace} la diagonalización a un
        sistema acoplado? ¿Y en qué sentido el estado estacionario de
        Markov de la pregunta 14 es también una historia de \omterm{def:g11:vect:vector}{vectores}
        propios?
  \item Final: las tres caras de la \omterm{def:g12:matrix:matrix}{matriz} este fin de semana: la
        contabilidad (sistemas e inversas), la combinatoria (caminos y
        enlaces contados por las potencias) y la evolución (Fibonacci, el
        tiempo, la web: futuros que se leen en las direcciones propias).
        Una frase para cada una, más la mirada hacia delante: el álgebra
        lineal de los volúmenes universitarios convierte cada una de esas
        caras en una teoría.
\end{enumerate}
\end{problem}
